% GENERATED FILE. Edit canonical lessons, then run npm run generate:books.
% canonical-source-hash: fnv1a64:c96df8f9ce7b36ba
% canonical-lessons: GE-C371-zusagen, GE-C371-zusage, GE-C371-absagen, GE-C371-absage, GE-C371-bestaetigung

\chapter{Zusagen und absagen}
\label{ch:ge-zusagen-und-absagen}
\hlchaptermodality{371}

\begin{chapteropening}
\textbf{By the end of this chapter:} \emph{I can accept or turn down an invitation or an offer, call something off, and talk about a confirmation.}

Complete the last lesson of chapter 371: I can accept or turn down an invitation or an offer, call something off, and talk about a confirmation.
\end{chapteropening}

\section[zusagen]{zusagen — to accept (an invitation), to promise; to appeal to}

\label{lesson:GE-C371-zusagen}



\begin{quote}
Before the new one: say the German for to be true, then the German for to turn out.
\end{quote}

\subsection*{You'll want to know: zusagen}
\textbf{zusagen} — \textquotedblleft{}to accept (an invitation), to promise; to appeal to\textquotedblright{}.

\begin{quote}
Sie hat für die Hochzeit zugesagt.
\end{quote}

\begin{quote}
\emph{She has accepted the invitation to the wedding.}
\end{quote}

Separable: \emph{sie sagt zu}, \emph{sie hat zugesagt}. \emph{Das sagt mir zu} — I like that.

\begin{cousinweb}
\emph{zusagen} is built from \emph{zu} and \emph{sagen} \textquotedblleft{}to say\textquotedblright{}, inherited from Proto-Germanic \emph{sagjaną} (EtymDB), a cousin of English \emph{say}: you say yes to something.
\end{cousinweb}

\begin{grammarlens}[title={What you've built}]
One more word for telling, explaining and describing at B1.
\end{grammarlens}

\subsection*{Your turn}
\begin{itemize}
\raggedright
  \item \emph{Say it:} \emph{zusagen}
  \item \emph{Say it:} \emph{zusagen}, once more
  \item \emph{Say it:} a sentence of your own with \emph{zusagen}
  \item \emph{Recall:} say the German for to be true, then the German for to turn out, then say \emph{zusagen} again
\end{itemize}

\begin{tcolorbox}[breakable,colback=teal!4,colframe=teal!35!black,title={Before you move on}]
What does zusagen mean? (\textquotedblleft{}To accept (an invitation), to promise; to appeal to\textquotedblright{}.) Say it once more.
\end{tcolorbox}
\section[die Zusage]{die Zusage — the acceptance, the promise; the job offer}

\label{lesson:GE-C371-zusage}



\begin{quote}
Before the new one: say the German for to turn out, then the German for to accept.
\end{quote}

\subsection*{You'll want to know: die Zusage}
\textbf{die Zusage} — \textquotedblleft{}the acceptance, the promise; the job offer\textquotedblright{}.

\begin{quote}
Nach dem Vorstellungsgespräch habe ich eine Zusage bekommen.
\end{quote}

\begin{quote}
\emph{After the interview I received a job offer.}
\end{quote}

Plural: \emph{die Zusagen}.

\begin{cousinweb}
\emph{Zusage} is the noun of \emph{zusagen}, from \emph{zu} and \emph{sagen} \textquotedblleft{}to say\textquotedblright{}, Proto-Germanic \emph{sagjaną} (EtymDB).
\end{cousinweb}

\begin{grammarlens}[title={What you've built}]
One more word for telling, explaining and describing at B1.
\end{grammarlens}

\subsection*{Your turn}
\begin{itemize}
\raggedright
  \item \emph{Say it:} \emph{die Zusage}
  \item \emph{Say it:} \emph{die Zusage}, once more
  \item \emph{Say it:} a sentence of your own with \emph{die Zusage}
  \item \emph{Recall:} say the German for to turn out, then the German for to accept, then say \emph{die Zusage} again
\end{itemize}

\begin{tcolorbox}[breakable,colback=teal!4,colframe=teal!35!black,title={Before you move on}]
What does die Zusage mean? (\textquotedblleft{}The acceptance, the promise; the job offer\textquotedblright{}.) Say it once more.
\end{tcolorbox}
\section[absagen]{absagen — to call off, to cancel; to decline}

\label{lesson:GE-C371-absagen}



\begin{quote}
Before the new one: say the German for to accept, then the German for the acceptance.
\end{quote}

\subsection*{You'll want to know: absagen}
\textbf{absagen} — \textquotedblleft{}to call off, to cancel; to decline\textquotedblright{}.

\begin{quote}
Wegen des schlechten Wetters wurde das Konzert abgesagt.
\end{quote}

\begin{quote}
\emph{Because of the bad weather the concert was called off.}
\end{quote}

Separable: \emph{er sagt ab}, \emph{er hat abgesagt}. \emph{Einen Termin absagen} — to cancel an appointment.

\begin{cousinweb}
\emph{absagen} is built from \emph{ab} \textquotedblleft{}off\textquotedblright{} and \emph{sagen} \textquotedblleft{}to say\textquotedblright{}, Proto-Germanic \emph{sagjaną} (EtymDB): you say a plan off.
\end{cousinweb}

\begin{grammarlens}[title={What you've built}]
One more word for telling, explaining and describing at B1.
\end{grammarlens}

\subsection*{Your turn}
\begin{itemize}
\raggedright
  \item \emph{Say it:} \emph{absagen}
  \item \emph{Say it:} \emph{absagen}, once more
  \item \emph{Say it:} a sentence of your own with \emph{absagen}
  \item \emph{Recall:} say the German for to accept, then the German for the acceptance, then say \emph{absagen} again
\end{itemize}

\begin{tcolorbox}[breakable,colback=teal!4,colframe=teal!35!black,title={Before you move on}]
What does absagen mean? (\textquotedblleft{}To call off, to cancel; to decline\textquotedblright{}.) Say it once more.
\end{tcolorbox}
\section[die Absage]{die Absage — the rejection; the cancellation}

\label{lesson:GE-C371-absage}



\begin{quote}
Before the new one: say the German for the acceptance, then the German for to call off.
\end{quote}

\subsection*{You'll want to know: die Absage}
\textbf{die Absage} — \textquotedblleft{}the rejection; the cancellation\textquotedblright{}.

\begin{quote}
Leider habe ich auf meine Bewerbung eine Absage bekommen.
\end{quote}

\begin{quote}
\emph{Unfortunately I got a rejection for my application.}
\end{quote}

Plural: \emph{die Absagen}. The opposite of \emph{die Zusage}.

\begin{cousinweb}
\emph{Absage} is the noun of \emph{absagen}, from \emph{ab} and \emph{sagen} \textquotedblleft{}to say\textquotedblright{}, Proto-Germanic \emph{sagjaną} (EtymDB).
\end{cousinweb}

\begin{grammarlens}[title={What you've built}]
One more word for telling, explaining and describing at B1.
\end{grammarlens}

\subsection*{Your turn}
\begin{itemize}
\raggedright
  \item \emph{Say it:} \emph{die Absage}
  \item \emph{Say it:} \emph{die Absage}, once more
  \item \emph{Say it:} a sentence of your own with \emph{die Absage}
  \item \emph{Recall:} say the German for the acceptance, then the German for to call off, then say \emph{die Absage} again
\end{itemize}

\begin{tcolorbox}[breakable,colback=teal!4,colframe=teal!35!black,title={Before you move on}]
What does die Absage mean? (\textquotedblleft{}The rejection; the cancellation\textquotedblright{}.) Say it once more.
\end{tcolorbox}
\section[die Bestätigung]{die Bestätigung — the confirmation}

\label{lesson:GE-C371-bestaetigung}



\begin{quote}
Before the new one: say the German for to call off, then the German for the rejection.
\end{quote}

\subsection*{You'll want to know: die Bestätigung}
\textbf{die Bestätigung} — \textquotedblleft{}the confirmation\textquotedblright{}.

\begin{quote}
Die Bestätigung für das Hotel kommt per E-Mail.
\end{quote}

\begin{quote}
\emph{The confirmation for the hotel comes by email.}
\end{quote}

Plural: \emph{die Bestätigungen}. From \emph{bestätigen} — to confirm.

\begin{cousinweb}
\emph{Bestätigung} is the noun of \emph{bestätigen} \textquotedblleft{}to confirm\textquotedblright{}, which is usually traced to the old adjective \emph{stet} \textquotedblleft{}steady, firm\textquotedblright{}, of the family of \emph{stehen} \textquotedblleft{}to stand\textquotedblright{}. EtymDB has no entries for these, so the account is hedged.
\end{cousinweb}

\begin{grammarlens}[title={What you've built}]
One more word for telling, explaining and describing at B1.
\end{grammarlens}

\subsection*{Your turn}
\begin{itemize}
\raggedright
  \item \emph{Say it:} \emph{die Bestätigung}
  \item \emph{Say it:} \emph{die Bestätigung}, once more
  \item \emph{Say it:} a sentence of your own with \emph{die Bestätigung}
  \item \emph{Recall:} say the German for to call off, then the German for the rejection, then say \emph{die Bestätigung} again
\end{itemize}

\begin{tcolorbox}[breakable,colback=teal!4,colframe=teal!35!black,title={Before you move on}]
What does die Bestätigung mean? (\textquotedblleft{}The confirmation\textquotedblright{}.) Say it once more.
\end{tcolorbox}
